\documentclass[11pt]{article}
\usepackage{geometry}
\geometry{verbose}
\setlength{\parskip}{6pt}
\setlength{\parindent}{0pt}
\usepackage{color}
\usepackage{amsmath, amssymb}
\usepackage[authoryear]{natbib}
\usepackage[unicode=true,
 bookmarks=false,
 breaklinks=false,pdfborder={0 0 1},backref=section,colorlinks=false]
 {hyperref}

\makeatletter
\newcommand{\R}{\mathbb{R}}
\newcommand{\E}{\mathbb{E}}
\providecommand{\tabularnewline}{\\}
\usepackage{graphicx}
\usepackage{float}
\usepackage{url}
\usepackage{array}
\usepackage{enumitem}

\title{\Huge Assignment 2}
\author{\Large Chris Conlon}
\date{\Large Fall 2026 --- due at the start of Session 4}

\makeatother

\begin{document}

\title{Assignment 2}
\maketitle
\begin{center}
\begin{tabular*}{0.9\textwidth}{@{\extracolsep{\fill}}@{\extracolsep{\fill}}l@{\extracolsep{\fill}}l@{\extracolsep{\fill}}l}
Econometrics I & $\qquad$ & Professor Chris Conlon\tabularnewline
NYU Stern &  & Email: \href{mailto:cconlon@stern.nyu.edu}{cconlon@stern.nyu.edu}\tabularnewline
\end{tabular*}
\par\end{center}

\textbf{Instructions.} This assignment covers Lecture 3 (inference and standard
errors). You may use \texttt{R} (with \texttt{fixest}) or Python (with
\texttt{pyfixest}). Turn in (1) a PDF presenting your results and showing your work,
and (2) your code.

\subsection*{1 Monte Carlo: sampling distributions, standard errors, and coverage}

Consider the data generating process
\[
y_{i} = 1 + 2 x_{i} + \varepsilon_{i}, \qquad x_{i} \sim \mathcal{N}\left(0,1\right)
\]
For each design below, simulate $S = 1{,}000$ datasets of size $n = 50$, estimate the
regression by OLS in each, and save $\hat{\beta}_{1}^{(s)}$ and its reported standard
error. Report: (i) the mean of $\hat{\beta}_{1}^{(s)}$; (ii) the standard deviation of
$\hat{\beta}_{1}^{(s)}$ across simulations (the ``true'' standard error); (iii) the
average reported standard error; and (iv) \emph{coverage}: the fraction of simulations
in which the 95\% confidence interval contains $\beta_{1} = 2$.
\begin{enumerate}[label=(\alph*)]
\item $\varepsilon_{i} \sim \mathcal{N}\left(0, 4\right)$, classical standard errors.
Does everything line up the way Lectures 1b and 3 promised?
\item Heteroskedasticity: $\varepsilon_{i} = x_{i}^{2} \cdot \nu_{i}$ with
$\nu_{i} \sim \mathcal{N}\left(0,1\right)$. Compute coverage using (i) classical and
(ii) heteroskedasticity-robust (HC1) standard errors. Which one is lying to you, and
in which direction?
\item Clustering: assign observations to $G = 25$ groups of 20 (so $n = 500$ here), and let
\[
\varepsilon_{ig} = \eta_{g} + \nu_{ig}, \qquad \eta_{g} \sim \mathcal{N}\left(0,1\right),
\quad \nu_{ig} \sim \mathcal{N}\left(0,1\right)
\]
with $x_{ig} = w_{g} + z_{ig}$ where $w_{g}, z_{ig} \sim \mathcal{N}(0,1)$ (so the
regressor is also correlated within group). Compute coverage using classical, HC1,
and cluster-robust standard errors. Present the three coverage rates in one table and
summarize the lesson in two sentences.
\item In (c), how does cluster-robust coverage change if you re-run the design with
$G = 5$ groups of 100? One sentence on what this says about ``number of clusters''
warnings you will encounter in applied work.
\end{enumerate}

\subsection*{2 Bootstrap vs. the formulas}

Return to design (a) of Question 1, but with a single simulated dataset of $n = 50$.
\begin{enumerate}[label=(\alph*)]
\item Compute the standard error of $\hat{\beta}_{1}$ three ways: the classical formula,
HC1, and a nonparametric bootstrap with $B = 2{,}000$ resamples (resample rows with
replacement). Compare.
\item Now compute all three for the \emph{ratio} $\hat{\beta}_{1} / \hat{\beta}_{0}$.
(For the formula-based versions, note there is no textbook formula --- explain in one
sentence why the bootstrap doesn't care.)
\end{enumerate}

\subsection*{3 Fuel bills: the CEF in the wild}

The file \texttt{fuelbills.csv} contains data on 144 homes: \texttt{FUELBILL} (annual
fuel expenditure) and \texttt{ROOMS} (number of rooms), among other variables.
\begin{enumerate}[label=(\alph*)]
\item Produce a scatterplot of \texttt{FUELBILL} against \texttt{ROOMS}. Compute and plot
the mean of \texttt{FUELBILL} for each value of \texttt{ROOMS}. (You have just estimated
a conditional expectation function.)
\item Regress \texttt{FUELBILL} on \texttt{ROOMS} (with \texttt{feols} or
\texttt{pf.feols}). Interpret the slope in one sentence. Overlay the fitted line on your
plot from (a): where does the linear approximation fit the CEF well, and where does it not?
\item Now run the \emph{saturated} regression of \texttt{FUELBILL} on a full set of
indicator variables for each value of \texttt{ROOMS} (in \texttt{fixest}:
\texttt{feols(FUELBILL \~{} i(ROOMS))}; in \texttt{pyfixest}:
\texttt{pf.feols("FUELBILL \~{} i(ROOMS)")}). Show that the fitted values reproduce the
conditional means from (a) exactly. Why must this be true?
\item Report the regression from (b) with classical and with heteroskedasticity-robust
standard errors. Is there visual evidence in (a) that the robust ones are the right choice?
\end{enumerate}

\end{document}
